\documentclass{conjura-conjecture}

\newcommand{\calK}{\mathcal{K}}
\newcommand{\calD}{\mathcal{D}}
\newcommand{\calR}{\mathcal{R}}
\newcommand{\Adv}{\mathbf{Adv}}
\newcommand{\sd}{\mathit{sd}}
\newcommand{\extpubk}{\textup{ext-pub-}k}

\runninghead{LHL EXTRACTION, K-SOURCE (TEST PLACEHOLDER)}

\begin{document}

\cjkicker{CONJURA \textperiodcentered{} OPEN PROBLEM (TEST PLACEHOLDER)}
\cjtitle{Leftover Hash Lemma Extraction Bound for Unpredictable Random-Oracle Sources, Public Seed, $k$-Source Generalization}
\cjstatus{Statement: AI-written, not yet formalized.}
\cjcategory{Information-Theoretic Cryptography.}

\vspace{0.4em}

\noindent\textbf{Placeholder.} This is a lattice test entry, not yet reviewed for provenance: it asks
whether joint multi-source extraction beats the naive $k$-fold hybrid bound obtained by repeating the
single-source ($k=1$) argument of
\href{https://crypto-conjura.github.io/c/0004/}{crypto-conjura.github.io/c/0004} $k$ times.

\section{Setting}

As in the $k=1$ base case: fix nonempty finite $\calK$ (seeds), $\calD$ (inputs), $\calR$ (outputs),
and write $K := |\calK|$, $D := |\calD|$, $R := |\calR|$. Take $k$ independent
$\epsilon$-unpredictable sources $S_1, \dots, S_k$ over the same $(\calK,\calD,\calR)$, each producing
its own $(x_i, z_i)$. Draw $k$ independent seeds $\sd_1, \dots, \sd_k \stackrel{\$}{\leftarrow} \calK$
and reveal all of them publicly, alongside all of $H$ and all the $z_i$.

\begin{definition}[$k$-Instance Public-Seed Extraction]
For a challenge bit $b$, set $y_{i,0} \leftarrow H(\sd_i, x_i)$ and
$y_{i,1} \stackrel{\$}{\leftarrow} \calR$ for each $i$, and let $\mathsf{D}$ output a guess $b'$ given
$H$, all $\sd_i$, all $y_{i,b}$, and all $z_i$. Define
\[
  \Adv^{\extpubk}_{\calK,\calD,\calR}(S_1,\dots,S_k,\mathsf{D}) := 2\Pr[b'=b] - 1 .
\]
\end{definition}

\section{The Placeholder Conjecture}

\begin{proposition}[Trivial hybrid bound, not in question]
By a standard hybrid argument over the $k$ instances, using the corollary of the $k=1$ base case at
each step,
\[
  \Adv^{\extpubk}_{\calK,\calD,\calR}(S_1,\dots,S_k,\mathsf{D})
  \;\le\; k \cdot \frac{8}{5}\Bigl( \sqrt{\epsilon R} + \sqrt{\tfrac{\log_2 D}{K}} \Bigr) .
\]
At $k=1$ this is exactly the base case's corollary, which is why that statement specializes this one.
\end{proposition}

\begin{conjecture}[Placeholder, untested]
The linear-in-$k$ factor above is not necessary: there is a bound of the same shape with the
$k$-dependence replaced by $O(\sqrt{k})$ or better, exploiting that a single distinguisher must
succeed jointly across all $k$ instances rather than being handed $k$ independent certificates.
\end{conjecture}

\section{Open Obligations}

\begin{itemize}
  \item Decide whether this is worth pursuing as a real statement (drop the test category and attach
  real provenance) or leave as a lattice placeholder.
  \item If pursued: attempt to beat the trivial $k$-fold hybrid bound, or show it is tight.
\end{itemize}

\end{document}
